\begin{scanpage}{0120}{103}
\[\sigma=\pi a_1^2\left|-ika_1-\frac1{n_1}\operatorname{ctg}\frac\pi{n_1}+\frac{\frac2{n_2}\sin\frac\pi{n_2}}{\cos\frac\pi{n_2}-1}e^{2ikl_2}\right|^2,\tag{17.25}\]
\[\sigma^0=\pi a_1^2(ka_1)^2,\tag{17.26}\]
in connection with which
\[n_1=\frac32,\quad n_2=1+\frac\Omega\pi.\tag{17.27}\]
Equation (17.25) is more precise than Equation (15.06) which was derived in \S\ 15, where the value of the field in the direction $\vartheta=\pi$ was taken in the physical optics approach.
\sectiontitle{18}{A Paraboloid of Rotation}
Let us calculate the effective scattering area of a paraboloid of rotation $r^2=2pz$ (Figure 36) which is irradiated by plane wave (16.20). The uniform part of the current excited on the paraboloid's surface has the components
\[\left.\begin{aligned}j_x^0&=\frac c{2\pi}E_{0x}\sin\alpha\,e^{ikz},\\j_y^0&=0,\\j_z^0&=\frac c{2\pi}E_{0x}\cos\alpha\cos\psi\,e^{ikz}.\end{aligned}\right\}\tag{18.01}\]
Integrating this current, it is not difficult to show that in the direction $\vartheta=\pi$ it radiates the field
\[\left.\begin{aligned}E_x=-H_y&=-E_{0x}\frac a2(1-e^{2ikl})\operatorname{tg}\omega\frac{e^{ikR}}R,\\E_y=H_x&=0.\end{aligned}\right\}\tag{18.02}\]
Here $a$ is the radius of the base of the paraboloid; $l=\frac{a^2}{2p}=\frac a2\operatorname{ctg}\omega$ is its length; $\alpha$ is the angle between the $z$ axis and the tangent to the generatrix of the paraboloid $(r^2=2pz)$. At the point $z=l$, the angle $\alpha$ takes the value $\alpha=\omega$ $(\operatorname{tg}\omega=\frac pa)$.
\end{scanpage}
